% ===== Wrap-up (~4 min, then questions) =====
\partaccent{butter}\notopic

\begin{frame}{Այսօրվա 4 գաղափարը}
  \begin{enumerate}
    \item<1-> Կարճաժամկետ հիշողությունը տեղավորում է \textbf{մոտ 4} բան, իսկ երկարաժամկետի սահմանը դեռ ոչ ոք չի գտել
    \item<2-> Դեմքը նկար է, անունը՝ ձայն. անունը պետք է \textbf{կրկնել}, թե չէ զրույցը այն դուրս է հրում
    \item<3-> Ալկոհոլը նախ անջատում է «կարմիր լույսերը», հետո՝ \textbf{հիշողության գրանցումը}
    \item<4-> Հիշողությունը ձայնագրություն չէ. այն ամեն անգամ \textbf{վերակառուցվում} ու մի քիչ խմբագրվում է, հաճախ՝ մեր օգտին
  \end{enumerate}
\end{frame}

\begin{frame}{Փորձիր տանը}
  \begin{columns}[T]
    \begin{column}{0.31\textwidth}
      \begin{card}[colback=lavender!35!cardwhite]
        \centering\bigemoji[1cm]{wavehand}\par\smallskip
        \raggedright Նոր ծանոթի անունը կրկնիր ու զրույցում երկու անգամ օգտագործիր։ Վաղը կհիշե՞ս
      \end{card}
    \end{column}
    \begin{column}{0.31\textwidth}
      \begin{card}[colback=mint!30!cardwhite]
        \centering\bigemoji[1cm]{numbers}\par\smallskip
        \raggedright Նոր հեռախոսահամարը հիշիր երեք խմբով, ոչ թե թիվ առ թիվ։ Ավելի հե՞շտ է
      \end{card}
    \end{column}
    \begin{column}{0.31\textwidth}
      \begin{card}[colback=peach!35!cardwhite]
        \centering\bigemoji[1cm]{family}\par\smallskip
        \raggedright Ընտանիքիդ հետ համեմատիր մի ընդհանուր հին հիշողություն։ Նույն կերպ եք հիշո՞ւմ
      \end{card}
    \end{column}
  \end{columns}
\end{frame}

\begin{frame}{Հաջորդ վեբինարը}
  \begin{columns}[c]
    \begin{column}{0.6\textwidth}
      \chip[brainpink]{scream}{3. Վախ. վախենալու բան չկա}\par\smallskip
      Սնահավատություն ու դավադրության տեսություններ, ֆոբիաներ, ինչու ենք սիրում վախենալ, ու ինչու է քննադատությունն ավելի ուժեղ, քան գովեստը\par\bigskip
      {\small\color{plum!70!cream}Այնուհետև ամեն վեբինարում՝ մեկ գլուխ. 4. Խելք, 5. Ընկալում, 6. Անհատականություն, 7. Ուրիշ մարդիկ, 8. Երբ ուղեղը «փչանում» է}
    \end{column}
    \begin{column}{0.34\textwidth}
      \centering\mascot[2.8cm]{scared}
    \end{column}
  \end{columns}
\end{frame}

\begin{frame}[plain]
  \begin{tikzpicture}[remember picture, overlay]
    \fill[lavender!60!cream] (current page.south west) rectangle (current page.north east);
    \fill[pattern={Dots[distance=4.4mm, radius=0.28mm]}, pattern color=plum!10!lavender!60!cream]
      (current page.south west) rectangle (current page.north east);
    \node at ([xshift=-3.2cm]current page.center) {\mascot[4.4cm]{happy}};
    \node[anchor=west, text=plum] at ([xshift=0.2cm, yshift=0.6cm]current page.center)
      {\titlefont\fontsize{40}{42}\selectfont Հարցեր};
    \node[anchor=west, text=plum, font=\mediumfont\large, align=left] at ([xshift=0.2cm, yshift=-1.0cm]current page.center)
      {Գրեք չաթում կամ\\ միացրեք խոսափողը \emoji{megaphone}};
    \node[anchor=west, text=plum!70!cream, font=\small] at ([xshift=0.2cm, yshift=-2.1cm]current page.center)
      {Դին Բըրնեթ, «The Idiot Brain», 2016};
  \end{tikzpicture}
\end{frame}
